\textbf{结论名称：} 等差数列片段和性质：\(S_m,\ S_{2m}-S_m,\ S_{3m}-S_{2m}\) 成等差数列。

\textbf{适用条件：} 数列 \(\{a_n\}\) 为等差数列，\(S_n\) 为其前 \(n\) 项和，\(m\in\mathbb{N}^*\)。

\textbf{变量说明：} 设原等差数列公差为 \(d\)；\(S_0=0\)；记片段和 \(T_k=S_{km}-S_{(k-1)m}\ (k=1,2,\dots)\)。

\textbf{核心公式：}
\[
\boxed{2(S_{2m}-S_m)=S_m+(S_{3m}-S_{2m})},\qquad 
\boxed{S_{(k+1)m}-S_{km}=S_m+k\cdot m^2 d\ \ (k\ge 0)}.
\]
片段和数列 \(\{T_k\}\) 是公差为 \(m^2 d\) 的等差数列。

\textbf{使用场景：} 已知等差数列前若干项和，快速求后续等间隔片段和；判断 \(S_{km}\) 型递推关系；处理“等距分段求和”的选择、填空题。

\textbf{简单例子：} 等差数列 \(\{a_n\}\) 中，\(a_1=1,\ d=2\)，取 \(m=3\)。则
\(S_3=1+3+5=9,\ S_6=1+\cdots+11=36,\ S_9=1+\cdots+17=81\)。
片段和：\(T_1=S_3=9,\ T_2=S_6-S_3=27,\ T_3=S_9-S_6=45\)。
\(27-9=18,\ 45-27=18\)，恰成公差 \(m^2 d=9\times2=18\) 的等差数列。
若仅知 \(\{a_n\}\) 为等差数列且 \(S_3=9,\ S_6=36\)，由性质直接得 \(T_3=2T_2-T_1=45\)，
从而 \(S_9=S_6+T_3=81\)。